\section{Scope and further directions}
\label{sec:scope}

The local criterion separates the four inputs that the proof actually uses.
A global upper density bound is unnecessary if root-test densities and root
incidences are controlled directly.  Thus variable-weighted or otherwise
nonuniform models may fall within the theorem even when their full clause
law is not bounded above by a constant multiple of the uniform law.  The
ordinary-clause lower density controls the potential drift, the test-clause
upper density controls pivotality, the incidence conditions control the
robustness factors and collisions, and the assignment-falsification bound is
needed only to remove the cap.

Under the simpler pointwise hypothesis, broad nonuniformity is allowed.  For
example, at time $j$ the unordered variable set may have probability between
fixed multiples of $\binom nk^{-1}$ and, conditional on that set, each sign
pattern may have probability between fixed multiples of $2^{-k}$.  No
convergence, periodicity, or stationarity of these laws is required.

Independence between clause times is used in the coordinate-omission
argument, the product structure of hidden root tests, the incidence moments,
and the uncapping tail.  Adaptive clause laws or an unconditioned common
random environment are therefore not covered as stated.

A matching linear lower variance bound remains the principal open question.
The mixture decomposition supplied by a pointwise lower density does not by
itself prevent the residual stream from determining the hitting time.  The
present proof also does not automatically yield subgaussian tails.  The
Boucheron--Bousquet--Lugosi--Massart moment method would require an almost
sure or sufficiently high-moment control of the full deletion variance
proxy, whereas here each $D_i$ can be of order $n$ and only
$\E D_i^2=O(1)$ is established.  Either strengthening would require a new
anti-concentration or high-moment ingredient rather than a formal reuse of
the variance proof.

\begingroup
\small
\bibliographystyle{plainnat}
\bibliography{references}
\endgroup
